\documentclass[aspectratio=169]{beamer}

% Tema
\usetheme{Madrid}
\usecolortheme{default}

% Paquetes
\usepackage[utf8]{inputenc}
\usepackage[spanish,es-tabla]{babel}
\usepackage{amsmath,amssymb,amsthm}
\usepackage{graphicx}
\usepackage{cleveref}
\usepackage{xcolor}

% Teoremas y definiciones
\newtheorem{teo}{Teorema}[section]
%\crefformat{teo}{\textbf{#2teorema~#1#3}} %Primera letra minuscula
\Crefformat{teo}{\textbf{#2Teorema~#1#3}}
\newtheorem{lem}{Lema}
\newtheorem{cor}{Corolario}
\theoremstyle{definition}
\newtheorem{defi}{Definición}
\newtheorem{eje}{Ejemplo}
\newtheorem{sol}{Solucion}

\bibliographystyle{plain}


% Información de la presentación
\title[Modelo Leslie-Gower con respuesta sigmoidea]{Un modelo presa-depredador tipo Leslie-Gower con respuesta funcional sigmoidea}
\author[Yerson Chilito]{Yerson Danilo Chilito Inga}
\institute[U. del Cauca]{
	Universidad del Cauca \\
	Departamento de Matemáticas \\
	\vspace{0.3cm}
	Director: Dr. Jhon Jairo Pérez
}
\date{2025}

\begin{document}
	
	% Portada
	\begin{frame}
		\titlepage
	\end{frame}
	
	% Índice
	\begin{frame}{Contenido}
		\tableofcontents
	\end{frame}
	
	% Introducción
	\section{Introducción}
	
	\begin{frame}{Introducción}
		\begin{block}{Contexto}
			Las ecuaciones diferenciales son instrumentos fundamentales para modelar fenómenos de la vida real, particularmente en ecología, física y procesos biológicos.
		\end{block}
		
		\vspace{0.5cm}
		
		\begin{itemize}
			\item La mayoría de estos modelos no admiten soluciones explícitas
			\item El análisis cualitativo es esencial para caracterizar el comportamiento dinámico
			\item Técnicas: localización de puntos de equilibrio, estabilidad, ciclos límite, bifurcaciones
		\end{itemize}
	\end{frame}
	
	\begin{frame}{El Modelo}
		\begin{block}{Sistema de Ecuaciones}
			\begin{equation*}
				\begin{cases} 
					\frac{dx}{dt} = \left({\color{red} r \left(1 - \frac{x}{K}\right)} - {\color{blue}\frac{q x y}{x^2 + a^2}} \right)x \\ 
					\frac{dy}{dt} = {\color{cyan} s \left(1 - \frac{y}{n x}\right)} y
				\end{cases}
			\end{equation*}
		\end{block}
		
		\textbf{Variables y Parámetros:}
		\begin{itemize}
			\item $x(t)$, $y(t)$: densidades de presa y depredador
			\item $r$: tasa intrínseca de crecimiento de la presa
			\item $K$: capacidad de carga ambiental
			\item $q$: tasa máxima de consumo por depredador
			\item $a$: cantidad de presas para media saturación
			\item $s$: tasa de crecimiento del depredador
			\item $n$: calidad alimenticia convertida en nuevos depredadores
		\end{itemize}
	\end{frame}
	
	\begin{frame}{Relevancia Ecológica}
		\begin{block}{Aplicaciones Reales}
			El modelo describe interacciones donde la respuesta funcional sigmoidea estabiliza las poblaciones.
		\end{block}
		
		\vspace{0.5cm}
		
		\textbf{Ejemplo:} Interacción entre focas y salmones en ríos
		\begin{itemize}
			\item Depredación mínima a bajas densidades de presa
			\item Depredación intensiva a densidades elevadas
		\end{itemize}
		
		\vspace{0.5cm}
		
		\begin{alertblock}{Importancia}
			Permite predecir la coexistencia o extinción de especies, ilustrando la aplicabilidad matemática en la conservación ecológica.
		\end{alertblock}
	\end{frame}
	
	% Marco Teórico
	\section{Marco Teórico}
	
	\begin{frame}{Ecuaciones Diferenciales Ordinarias}
		\begin{defi}[EDO]
			Una EDO es una ecuación para una función desconocida de una variable real, que contiene la función y sus derivadas: $F(t, x, x^{(1)}, \ldots, x^{(k)}) = 0$
		\end{defi}
		
		\vspace{0.3cm}
		
		\begin{defi}[Sistema Autónomo]
			Sistema de primer orden de la forma:
			\begin{equation}
				\dot{x} = f(x)
				\label{eq: sistema}
			\end{equation}
		\end{defi}
		
		\vspace{0.3cm}
		
		\begin{teo}[Existencia y Unicidad]
			Si $E \subseteq \mathbb{R}^n$ es abierto y $f \in C^1(E)$, entonces para cada $x_0 \in E$, el PVI tiene una única solución en algún intervalo.
		\end{teo}
	\end{frame}
	
	\begin{frame}{Conceptos Clave}
		\begin{columns}
			\column{0.5\textwidth}
			\begin{itemize}
				\item \textbf{Flujo y Trayectorias}: Conjunto de soluciones $\phi_t(x_0)$
				\item \textbf{Retrato Fase}: Conjunto de todas las curvas solución en $\mathbb{R}^n$
				\item \textbf{Isoclinas}: Curvas donde la pendiente es constante
				\item \textbf{Ciclos Límite}: Trayectorias cerradas y aisladas (estables, inestables o medio-estables)
			\end{itemize}
			\column{0.5\textwidth}
			\begin{figure}
				\centering
				\includegraphics[width=1\textwidth]{/home/skrap127/Documentos/1_/img/Example.pdf}
				\caption{Ejemplo de campo vectorial, trayectorias, retrato fase e isoclinas del sistema \eqref{eq: sistema} asociado a $f(x)=(x_{1},\sin{x_{2}})$.}
				\label{fig: example}
			\end{figure}
		\end{columns}
	\end{frame}
	
	\begin{frame}
		\begin{figure}
			\centering
			\includegraphics[width=0.6\textwidth]{/home/skrap127/Documentos/1_/img/CicloLimEstInest.pdf}
			\caption{Ejemplo de ciclos limites del sistema \eqref{eq: sistema} asociado a $f(x) = (-x_{2} + x_{1}(r^{4} -3r^{2} + 1), x_{1} + x_{2}(r^{4} -3r^{2} + 1))$ con $r^{2} = x_{1}^{2} + x_{2}^{2}$.}
			\label{fig: CicloLimiteEstableInestable}
		\end{figure}
	\end{frame}
	
	\begin{frame}{Puntos de Equilibrio}
		\begin{defi}[Punto de Equilibrio]
			$x_0 \in \mathbb{R}^n$ es un punto de equilibrio de \eqref{eq: sistema}, si $f(x_0) = 0$. Además, se le llama hiperbólico si ninguno de los valores propios de la matriz jacobiana $Df(x_0)$ tiene parte real nula, en caso contrario se le dice no hiperbólico.
		\end{defi}
		
		\vspace{0.3cm}
		
		\textbf{Clasificación (puntos hiperbólicos):}
		Depende de los valores propios de la matriz jacobiana $Df(x_0)$.
		\begin{itemize}
			\item \textbf{Sumidero}: Todos los valores propios con parte real negativa
			\item \textbf{Fuente}: Todos los valores propios con parte real positiva
			\item \textbf{Punto Silla}: Valores propios con partes reales de signos opuestos
		\end{itemize}
		
		Otros conceptos relacionados son:
		
		\begin{itemize}
			\item \textbf{Variedades Estables/Inestables}: Describen puntos que tienden al equilibrio cuando $t \to \infty$ o $t \to -\infty$
			\item \textbf{Separatrices}: Dividen el plano en regiones con diferentes comportamientos
		\end{itemize}
	\end{frame}
	
	\begin{frame}
		\begin{figure}
			\centering
			\includegraphics[width=1\textwidth]{/home/skrap127/Documentos/1_/img/EquilibriumPoints.pdf}
			\caption{Ejemplo de puntos de equilibrio en el origen de sistemas lineales bidimensionales.}
			\label{fig: EquilibriumPoints}
		\end{figure}
	\end{frame}
	
	\begin{frame}{Teoremas Fundamentales}
		\begin{teo}[Poincaré-Bendixson]
			Si una trayectoria está confinada en un subconjunto cerrado y acotado del plano sin puntos de equilibrio, entonces contiene una órbita cerrada.
		\end{teo}
		
		\vspace{0.5cm}
		
		\begin{teo}[Hartman-Grobman]
			El retrato fase cerca de un punto hiperbólico es topológicamente equivalente al de su linealización.
		\end{teo}
	\end{frame}
	
	\begin{frame}
		\begin{figure}
			\centering
			\includegraphics[width=1\textwidth]{/home/skrap127/Documentos/1_/img/exampleHartman.pdf}
			\caption{Retrato fase del sistema \eqref{eq: sistema} asociado a $f(x,y) = (-x, y+x^2)$ y su respectiva transformación en el origen dada por el homeomorfismo $H(x,y)=(x,x+\frac{y^{2}}{3})$}
			\label{fig: exampleHartman}
		\end{figure}
	\end{frame}
	
	\begin{frame}{Técnicas Avanzadas}
		\begin{block}{Blow-up (Desingularización)}
			Herramienta para estudiar singularidades complicadas mediante transformación del punto crítico en un círculo con nuevos puntos críticos (a menudo hiperbólicos).
		\end{block}
		
		\vspace{0.5cm}
		
		\begin{block}{Compactificación de Poincaré}
			Método que compactifica el espacio cartesiano extendiendo el campo vectorial a una esfera unitaria $S^2$, permitiendo analizar el comportamiento cuando las variables tienden a infinito.
		\end{block}
	\end{frame}
	
	\begin{frame}
		\begin{figure}
			\centering
			\includegraphics[width=.6\textwidth]{/home/skrap127/Documentos/1_/img/poincareCompactification.pdf}
			\caption{Esfera de Poincaré.}
			\label{fig: poincareCompactification}
		\end{figure}
	\end{frame}
	
	\begin{frame}{Bifurcaciones}
		\begin{defi}[Bifurcación]
			Cambio en la estructura cualitativa del flujo del sistema \begin{equation}
				\dot{x} = f_{\mu}(x) \quad x \in \mathbb{R}^{n}, \quad \mu \in \mathbb{R}^{k}
				\label{eq: sisParametro}
			\end{equation}
			cuando un parámetro $\mu$ varía.
		\end{defi}
		
		\vspace{0.3cm}
		
		\textbf{Tipos principales:}
		\begin{itemize}
			\item \textbf{Silla-Nodo}: Creación y destrucción de puntos de equilibrio
			\item \textbf{Transcrítica}: Cambio de estabilidad de punto fijo existente
			\item \textbf{Pitchfork}: Aparición de pares simétricos (supercrítica/subcrítica)
			\item \textbf{Hopf}: Punto fijo estable se vuelve inestable, aparecen ciclos límite (supercrítica/subcrítica)
		\end{itemize}
	\end{frame}
	
	\begin{frame}
		\begin{figure}[H]
			\centering
			\includegraphics[width=1\textwidth]{/home/skrap127/Documentos/1_/img/bifHopfSup.pdf}
			\caption{Ejemplo de bifurcación de Hopf supercrítica en el sistema \eqref{eq: sisParametro} asociado a $f_{\mu}(r,\theta)=(\mu r - r^3,1 + r^2)$.}
			\label{fig: bifHopfSup}
		\end{figure}
	\end{frame}
	
	\begin{frame}{Modelos Ecológicos: Crecimiento Poblacional}
		\begin{itemize}
			\item \textbf{Crecimiento Exponencial}: Poco realista a largo plazo, ignora limitaciones ambientales
			\item \textbf{Crecimiento Logístico}: Más relevante, el crecimiento es limitado por la capacidad de carga del entorno
		\end{itemize}
	\end{frame}
	
	\begin{frame}{Respuestas Funcionales de Holling}
		\begin{defi}[Respuesta Funcional]
			Relación entre la tasa de consumo de presas y el beneficio para el depredador.
		\end{defi}
		
		\vspace{0.3cm}
		
		\textbf{Tres tipos principales:}
		\begin{itemize}
			\item \textbf{Tipo I}: Aumento lineal con la densidad de presas
			\item \textbf{Tipo II}: Aumento con desaceleración gradual hasta una meseta
			\item \textbf{Tipo III (Sigmoidea)}: Fase de aceleración a bajas densidades (forma de S), luego meseta. El depredador aumenta su eficiencia cuando la presa es abundante
		\end{itemize}
	\end{frame}
	
	\begin{frame}
		\begin{columns}
			\column{0.5\textwidth}
			\begin{figure}[H]
				\centering
				\includegraphics[width=.5\textwidth]{/home/skrap127/Documentos/1_/img/type1FunctionalResponse.pdf}
				\caption{Respuesta funcional de Holling tipo I.}
				\label{fig: type1FunctionalResponse}
			\end{figure}
			\begin{figure}[H]
				\centering
				\includegraphics[width=.5\textwidth]{/home/skrap127/Documentos/1_/img/type2FunctionalResponse.pdf}
				\caption{Respuesta funcional de Holling tipo II.}
				\label{fig: type2FunctionalResponse}
			\end{figure}
			\column{0.5\textwidth}
			\begin{figure}[H]
				\centering
				\includegraphics[width=1\textwidth]{/home/skrap127/Documentos/1_/img/type3FunctionalResponse.pdf}
				\caption{Respuesta funcional de Holling tipo III.}
				\label{fig: type3FunctionalResponse}
			\end{figure}
		\end{columns}
	\end{frame}
	
	\begin{frame}{Modelos Clásicos Presa-Depredador}
		\begin{block}{Modelo Lotka-Volterra}
			\begin{equation*}
				\begin{cases} 
					\frac{dN}{dt} = N(a - bP) \\
					\frac{dP}{dt} = P(cN - d)
				\end{cases}
			\end{equation*}
			Crecimiento maltusiano de la presa, sin limitación de recursos.
		\end{block}
		
		\vspace{0.3cm}
		
		\begin{block}{Modelo Leslie-Gower}
			Variante que incorpora crecimiento logístico y capacidad de carga ($K$) del entorno para las presas, y capacidad de carga para depredadores dependiente de las presas ($K_P = dN$).
		\end{block}
	\end{frame}
	
	% Objetivos
	\section{Objetivos}
	
	\begin{frame}{Objetivo General}
		\begin{block}{}
			\centering
			\large
			Analizar cualitativamente un modelo presa-depredador tipo Leslie-Gower con respuesta funcional sigmoidea.
		\end{block}
	\end{frame}
	
	\begin{frame}{Objetivos Específicos}
		\begin{enumerate}
			\item Hallar puntos críticos del sistema
			\vspace{0.3cm}
			\item Analizar la estabilidad de dichos puntos críticos
			\vspace{0.3cm}
			\item Analizar la posibilidad de hallar comportamientos cualitativos especiales en el sistema (ciclos límites, bifurcaciones, etc.)
			\vspace{0.3cm}
			\item Determinar el comportamiento a futuro de las especies
		\end{enumerate}
	\end{frame}
	
	% Metodología
	\section{Metodología}
	
	\begin{frame}{Metodología}
		\begin{block}{Tipo de Investigación}
			Investigación documental que se fundamentará en el análisis matemático del modelo propuesto por González-Olivares et al. (2015).
		\end{block}
		
		\vspace{0.5cm}
		
		\textbf{Etapas del trabajo:}
		\begin{enumerate}
			\item Revisión bibliográfica
			\item Análisis y síntesis del marco teórico
			\item Aplicación de técnicas de análisis cualitativo
			\item Elaboración del documento final
		\end{enumerate}
	\end{frame}
	
	% Recursos y Presupuesto
	\section{Recursos y Presupuesto}
	
	\begin{frame}{Recursos}
		\begin{columns}
			\column{0.5\textwidth}
			\textbf{Recursos Humanos:}
			\begin{itemize}
				\item Director: Dr. Jhon Jairo Pérez
				\item Responsable: Yerson Danilo Chilito Inga
			\end{itemize}
			
			\column{0.5\textwidth}
			\textbf{Recursos Materiales:}
			\begin{itemize}
				\item Computador personal
				\item Libros de biblioteca
				\item Textos digitales
				\item Software especializado
			\end{itemize}
		\end{columns}
		
		\vspace{0.8cm}
		
		\begin{block}{Financiación}
			A cargo del estudiante. La Universidad del Cauca proporciona salas de estudio, bibliotecas y espacios básicos. El Departamento de Matemáticas asigna labor académica al director.
		\end{block}
	\end{frame}
	
	\begin{frame}{Cronograma}
		\begin{table}
			\centering
			\small
			\begin{tabular}{|l|c|c|c|c|c|c|}
				\hline
				\textbf{Actividad} & \textbf{M1} & \textbf{M2} & \textbf{M3} & \textbf{M4} & \textbf{M5} & \textbf{M6} \\ \hline
				Revisión bibliográfica & X & X & & & & \\ \hline
				Análisis y síntesis & & & X & X & & \\ \hline
				Documento final & & & & & X & X \\ \hline
				Sustentación & & & & & & X \\ \hline
			\end{tabular}
		\end{table}
	\end{frame}
	
	% Referencias
	\section{Referencias}
	
	\begin{frame}{Referencias Principales}
		\footnotesize
		\begin{itemize}
			\item González-Olivares, E., Tintinago-Ruiz, P. C., y Rojas-Palma, A. (2015). A Leslie–Gower-type predator–prey model with sigmoid functional response. \textit{International Journal of Computer Mathematics}, 92(9), 1895--1909.
			\item Murray, J. D. (2002). \textit{Mathematical Biology I: An Introduction} (3ra ed.). Springer.
			\item Perko, L. (2001). \textit{Differential Equations and Dynamical Systems} (3ra ed.). Springer.
			\item Strogatz, S. H. (2014). \textit{Nonlinear Dynamics and Chaos}. Westview Press.
			\item Teschl, G. (2012). \textit{Ordinary Differential Equations and Dynamical Systems}. Springer.			
		\end{itemize}
	\end{frame}
	
	% Cierre
	\begin{frame}[plain]
		\begin{center}
			\Huge ¡Gracias por su atención!
			
			\vspace{1cm}
			
			\Large Preguntas
		\end{center}
	\end{frame}
	
\end{document}